% problem_id: S001-groupoids-lemma-etale
% source_id: S001 (Stacks Project)
% source_file: groupoids.tex
% statement_locator: groupoids.tex lines 4568--4600
% proof_locator: groupoids.tex lines 4602--4742
% env: lemma
% id_note: label 跨文件重名 ⇒ id 用文件限定形式
% pairing: tight (verified)
% status: complete (statement + author proof verbatim + reason + locators)
%
\begin{lemma}
\label{lemma-etale}
Let $S$ be a scheme. Let $f : (U', R', s', t') \to (U, R, s, t, c)$ be a
morphism of groupoid schemes over $S$.
\begin{enumerate}
\item $U$, $R$, $U'$, $R'$ are affine,
\item $s, t, s', t'$ are finite locally free,
\item the diagrams
$$
\xymatrix{
R' \ar[d]_{s'} \ar[r]_f & R \ar[d]^s \\
U' \ar[r]^f & U
}
\quad
\quad
\xymatrix{
R' \ar[d]_{t'} \ar[r]_f & R \ar[d]^t \\
U' \ar[r]^f & U
}
\quad
\quad
\xymatrix{
G' \ar[d] \ar[r]_f & G \ar[d] \\
U' \ar[r]^f & U
}
$$
are cartesian where $G$ and $G'$ are the stabilizer group schemes, and
\item $f : U' \to U$ is \'etale.
\end{enumerate}
Then the map $C \to C'$ from the $R$-invariant functions on $U$
to the $R'$-invariant functions on $U'$ is \'etale and
$U' = \Spec(C') \times_{\Spec(C)} U$.
\end{lemma}

\begin{proof}
Set $M = \Spec(C)$ and $M' = \Spec(C')$.
Write $U = \Spec(A)$, $U' = \Spec(A')$, $R = \Spec(B)$, and
$R' = \Spec(B')$. We will use the results of
Lemmas \ref{lemma-integral-over-invariants},
\ref{lemma-invariants-base-change}, and
\ref{lemma-points}
without further mention.

\medskip\noindent
Assume $C$ is a strictly henselian local ring. Let $p \in M$
be the closed point and let $p' \in M'$ map to $p$.
Claim: in this case there is a disjoint union decomposition
$(U', R', s', t', c') = (U, R, s, t, c) \amalg (U'', R'', s'', t'', c'')$
over $(U, R, s, t, c)$ such that for the corresponding
disjoint union decomposition $M' = M \amalg M''$ over $M$
the point $p'$ corresponds to $p \in M$.

\medskip\noindent
The claim implies the lemma. Suppose that $M_1 \to M$ is a flat morphism
of affine schemes. Then we can base change everything to $M_1$
without affecting the hypotheses (1) -- (4).
From Lemma \ref{lemma-invariants-base-change}
we see $M_1$, resp.\ $M_1'$ is the spectrum of the
$R_1$-invariant functions on $U_1$,
resp.\ the $R'_1$-invariant functions on $U'_1$.
Suppose that $p' \in M'$ maps to $p \in M$.
Let $M_1$ be the spectrum of the strict henselization of
$\mathcal{O}_{M, p}$ with closed point $p_1 \in M_1$.
Choose a point $p'_1 \in M'_1$ mapping to $p_1$ and $p'$.
From the claim we get
$$
(U'_1, R'_1, s'_1, t'_1, c'_1) =
(U_1, R_1, s_1, t_1, c_1) \amalg
(U''_1, R''_1, s''_1, t''_1, c''_1)
$$
and correspondingly $M'_1 = M_1 \amalg M''_1$ as a scheme over $M_1$.
Write $M_1 = \Spec(C_1)$ and write $C_1 = \colim C_i$ as a filtered
colimit of \'etale $C$-algebras. Set $M_i = \Spec(C_i)$.
The $M_1 = \lim M_i$ and similarly for the other schemes.
By Limits, Lemmas \ref{limits-lemma-descend-opens} and
\ref{limits-lemma-descend-isomorphism}
we can find an $i$ such that
$$
(U'_i, R'_i, s'_i, t'_i, c'_i) =
(U_i, R_i, s_i, t_i, c_i) \amalg
(U''_i, R''_i, s''_i, t''_i, c''_i)
$$
We conclude that $M'_i = M_i \amalg M''_i$. In particular
$M' \to M$ becomes \'etale at a point over $p'$ after an
\'etale base change. This implies that $M' \to M$ is \'etale at $p'$
(for example by Morphisms, Lemma
\ref{morphisms-lemma-set-points-where-fibres-etale}).
We will prove $U' \cong M' \times_M U$ after we prove the claim.

\medskip\noindent
Proof of the claim. Observe that $U_p$ and $U'_{p'}$ have finitely many points.
For $u \in U_p$ we have $\kappa(u)/\kappa(p)$ is algebraic,
hence $\kappa(u)$ is separably closed.
As $U' \to U$ is \'etale, we conclude the morphism $U'_{p'} \to U_p$
induces isomorphisms on residue field extensions.
Let $u' \in U'_{p'}$ with image $u \in U_p$.
By assumption (3) the morphism of scheme theoretic fibres
$(s')^{-1}(u') \to s^{-1}(u)$,
$(t')^{-1}(u') \to t^{-1}(u)$, and
$G'_{u'} \to G_u$ are isomorphisms. Observing that $U_p = t(s^{-1}(u))$
(set theoretically) we conclude that the points of $U'_{p'}$
surject onto the points of $U_p$.
Suppose that $u'_1$ and $u'_2$ are points of $U'_{p'}$ mapping
to the same point $u$ of $U_p$. Then there exists a point
$r' \in R'_{p'}$ with $s'(r') = u'_1$ and $t'(r') = u'_2$.
Consider the two towers of fields
$$
\kappa(r')/\kappa(u'_1)/\kappa(u)/\kappa(p) \quad
\kappa(r')/\kappa(u'_2)/\kappa(u)/\kappa(p)
$$
whose ``ends'' are the same as the two ``ends'' of the two towers
$$
\kappa(r')/\kappa(u'_1)/\kappa(p')/\kappa(p) \quad
\kappa(r')/\kappa(u'_2)/\kappa(p')/\kappa(p)
$$
These two induce the same maps $\kappa(p') \to \kappa(r')$ as
$(U'_{p'}, R'_{p'}, s', t', c')$ is a groupoid over $p'$.
Since $\kappa(u)/\kappa(p)$ is purely inseparable,
we conclude that the two induced maps
$\kappa(u) \to \kappa(r')$ are the same.
Therefore $r'$ maps to a point of the fibre $G_u$.
By assumption (3) we conclude that $r' \in (G')_{u'_1}$.
Namely, we may think of $G$ as a closed subscheme of $R$
viewed as a scheme over $U$ via $s$ and use that
the base change to $U'$ gives $G' \subset R'$.
In particular we have $u'_1 = u'_2$.
We conclude that $U'_{p'} \to U_p$ is a bijective
map on points inducing isomorphisms on residue fields.
It follows that $U'_{p'}$ is a finite set of closed points
(Algebra, Lemma \ref{algebra-lemma-finite-residue-extension-closed})
and hence $U'_{p'}$ is closed in $U'$.
Let $J' \subset A'$ be the radical ideal cutting out $U'_{p'}$
set theoretically.

\medskip\noindent
Second part proof of the claim.
Let $\mathfrak m \subset C$ be the maximal ideal.
Observe that $(A, \mathfrak m A)$ is a henselian pair by
More on Algebra, Lemma \ref{more-algebra-lemma-integral-over-henselian-pair}.
Let $J = \sqrt{\mathfrak m A}$.
Then $(A, J)$ is a henselian pair
(More on Algebra, Lemma \ref{more-algebra-lemma-change-ideal-henselian-pair})
and the \'etale ring map
$A \to A'$ induces an isomorphism $A/J \to A'/J'$
by our deliberations above.
We conclude that $A' = A \times A''$ by
More on Algebra, Lemma \ref{more-algebra-lemma-characterize-henselian-pair}.
Consider the corresponding disjoint union
decomposition $U' = U \amalg U''$. The open $(s')^{-1}(U)$ is the
set of points of $R'$ specializing to a point of $R'_{p'}$.
Similarly for $(t')^{-1}(U)$. Similarly we have
$(s')^{-1}(U'') = (t')^{-1}(U'')$ as this is the set of
points which do not specialize to $R'_{p'}$.
Hence we obtain a disjoint union decomposition
$$
(U', R', s', t', c') =
(U, R, s, t, c) \amalg
(U'', R'', s'', t'', c'')
$$
This immediately gives $M' = M \amalg M''$ and the proof of the claim
is complete.

\medskip\noindent
We still have to prove that the canonical map $U' \to M' \times_M U$
is an isomorphism. It is an \'etale morphism
(Morphisms, Lemma \ref{morphisms-lemma-etale-permanence}).
On the other hand, by base changing to strictly henselian local rings
(as in the third paragraph of the proof) and using the bijectivity
$U'_{p'} \to U_p$ established in the course of the proof of the claim,
we see that $U' \to M' \times_M U$ is universally bijective
(some details omitted). However, a universally bijective
\'etale morphism is an isomorphism
(Descent, Lemma \ref{descent-lemma-universally-injective-etale-open-immersion})
and the proof is complete.
\end{proof}
